\chapterrule
\section{现代 Transformer 为什么看起来越来越不像原版？}
\label{intro:22d-modern-transformers}
\chaptercompass{现代架构替换了哪些 Transformer 部件？}{用明确的头数、缓存和计算形状比较 MHA、GQA、FFN 等设计。}{$Q:T\times H_Q\times d_k,\quad KV:T\times H_{KV}\times d_k$}{Q/K/V 投影、Norm 与 FFN 参数由具体配置确定}
\asof{2026-10-25}{相对位置、归一化与缓存预算}


先按顺序完成 RoPE、RMSNorm、GQA 和三维数值复核，再根据瓶颈选择后面的扩展专题。\par\smallskip


\subsection*{基础必读：位置、归一化与 KV 缓存}
\subsection{必读 1：RoPE 怎样利用相对位置？}
\phantomsection\label{concept:rope}

经典 Transformer 可以把位置向量 $P_{\rm pos}$ 加在 Token 嵌入上：
$X_0=X_{\rm token}+P_{\rm pos}$。现在进一步问：同样的两个词整体后移，
它们之间仍只隔一位，能否让注意力中的位置作用保持相同？
我们希望匹配能利用\textbf{相对距离}，而不只接收两个绝对位置编号。

一种办法是让位置决定向量的旋转角度：两个位置一起平移时，两向量多转相同角度，夹角不变。
这就是\textbf{旋转位置表示（RoPE）}的切入点。
先只看相同内容向量 $(1,0)$，它旋转 $\phi$ 后为 $(\cos\phi,\sin\phi)$。

\begin{example}
\noindent\emph{平移两个位置，分数不变}\par

取 $\vect q^\top=\vect k^\top=(1,0)$，角频率 $\theta=\pi/3$。
位置从 $1$ 起计，先放 Query 于 $m=1$、Key 于 $n=2$，分别旋转 $60^\circ$ 和 $120^\circ$：
\[
\bigl(\tfrac12,\tfrac{\sqrt3}2\bigr)\cdot
\bigl(-\tfrac12,\tfrac{\sqrt3}2\bigr)=-\tfrac14+\tfrac34=\boxed{\tfrac12}.
\]
再把两者整体后移两位，$m=3$、$n=4$，分别旋转 $180^\circ$ 和 $240^\circ$：
\[
(-1,0)\cdot\bigl(-\tfrac12,-\tfrac{\sqrt3}2\bigr)=\boxed{\tfrac12}.
\]
绝对位置变了，分数相同，因为两者都等于 $\cos\bigl((n-m)\theta\bigr)=\cos(\pi/3)$。
距离改成 $2$（如 $m=1,n=3$）时分数变为 $\cos(2\pi/3)=-\tfrac12$：
\textbf{分数只随相对距离变化}，而加在输入上的绝对位置向量一般不具备这种平移不变性。
\end{example}

图~\ref{fig:rope-relative-rotation} 把上面两次计算画在一起：
内容向量、角频率与位置取值相同，两次点积都为 $1/2$。
% Same unit vectors, angular frequency and position pairs as the preceding prose.
\par\noindent\begin{minipage}{\linewidth}
\tagstructbegin{tag=Figure,alttext={沿用正文的二维 RoPE 算例，内容向量均为一、零，角频率为三分之派。左图 Query 在位置零，旋转后为一、零；Key 在位置一，旋转后为二分之一、二分之根号三。右图两位置同时后移两位，Query 在位置二，旋转后为负二分之一、二分之根号三；Key 在位置三，旋转后为负一、零。两图的相对夹角均为六十度，相对位置差均为一，点积均为二分之一。}}
\begin{center}
\begin{tikzpicture}[>=Latex,x=14mm,y=14mm]
\foreach \shift/\title in {0/{位置 $m=0,\ n=1$},8.5/{位置 $m=2,\ n=3$}}{
 \begin{scope}[xshift=\shift cm]
 \draw[->,linegray] (-1.5,0)--(1.65,0);
 \draw[->,linegray] (0,-1.2)--(0,1.7);
 \draw[linegray,dashed] (0,0) circle (1);
 \node[font=\small\sffamily\bfseries] at (0,2.05) {\title};
 \node[font=\small] at (0,-1.6) {$n-m=1$，点积 $1/2$};
 \end{scope}
}
\begin{scope}
 \draw[->,very thick,blue] (0,0)--(0:1) node[below right,font=\small] {Q $(1,0)$};
 \draw[->,very thick,teal] (0,0)--(60:1) node[above,font=\small] {K $(\tfrac12,\tfrac{\sqrt3}2)$};
 \draw[->,flowcolor,thick] (.4,0) arc (0:60:.4);
 \node[font=\small] at (.65,.3) {$60^\circ$};
\end{scope}
\begin{scope}[xshift=8.5cm]
 \draw[->,very thick,blue] (0,0)--(120:1) node[above,font=\small] {Q $(-\tfrac12,\tfrac{\sqrt3}2)$};
 \draw[->,very thick,teal] (0,0)--(180:1) node[below left,font=\small] {K $(-1,0)$};
 \draw[->,flowcolor,thick] (120:.4) arc (120:180:.4);
 \node[font=\small] at (-.65,.3) {$60^\circ$};
\end{scope}
\node[font=\small,text=muted] at (3.0357,-2.35)
 {共同后移两位：两向量都多转 $120^\circ$，相对夹角仍为 $60^\circ$};
\end{tikzpicture}
\captionof{figure}{沿用正文的 $\theta=\pi/3$：$(m,n)$ 从 $(0,1)$ 平移到 $(2,3)$，点积仍为 $1/2$}
\label{fig:rope-relative-rotation}
\end{center}
\tagstructend
\end{minipage}\par


对一般内容向量，用二维旋转矩阵把刚才的操作写全：
\[
R(\phi)=\begin{bmatrix}\cos\phi&-\sin\phi\\\sin\phi&\cos\phi\end{bmatrix}.
\]
沿用行向量约定，在右侧乘旋转矩阵的转置。Query 与 Key 分别旋转 $m\theta$ 和 $n\theta$，则
\begin{equation}
(\vect q^\top R(m\theta)^\top)
(\vect k^\top R(n\theta)^\top)^\top
=\vect q^\top R((n-m)\theta)\vect k.
\label{eq:rope-relative}
\end{equation}
右边的位置项只含相对位置差 $n-m$；内容仍由 $q,k$ 决定。
高维 RoPE 对多个特征对使用不同角频率，把多种距离尺度一起带入匹配。

\subsection{必读 2：RMSNorm 为什么不减均值？}
\phantomsection\label{concept:rmsnorm}

\textbf{LayerNorm} 做了两件事：减均值，把中心移到零；再除以标准差，调整尺度。
如果我们主要想控制数值大小，是否一定要先把中心移到零？
可以直接用一行数的\textbf{均方根（RMS）}作为尺度：先逐项平方、取平均，再开平方。

只看一条二维向量 $\vect x=(3,4)$：
\[
\operatorname{RMS}(\vect x)
=\sqrt{\frac{3^2+4^2}{2}}=\sqrt{12.5}.
\]
忽略可学习缩放参数与小正数 $\varepsilon$，
相除后约为 $(0.849,1.131)$。若输入整体放大十倍为 $(30,40)$，
均方根也放大十倍，变成 $\sqrt{1250}$，相除仍得到 $(0.849,1.131)$。
因此，\textbf{直接除以均方根就能消去正的整体倍数，无需先减均值}。
这种归一化叫作\textbf{RMSNorm}，实际仍使用小正数 $\varepsilon$ 和可学习的逐特征缩放。

两种归一化保留的信息不同：对 $(3,4)$，忽略上述附加项的 LayerNorm 得到 $(-1,1)$，
RMSNorm 的两项则仍为正。RMSNorm 省去中心化、简化计算；
选择它的设计思路，是保留尺度控制，同时检验是否需要中心化。
具体模型的训练效果和运行成本要通过对照实验判断，原论文见\hyperref[source:rmsnorm]{附录来源}。


\textbf{Pre-Norm 与 Post-Norm} 的差别是归一化位于子层之前还是之后。
它改变残差与梯度路径，训练中的梯度因而经过不同的归一化位置。

\subsection{必读 3：GQA 怎样减少 KV 缓存？}
\phantomsection\label{concept:mla}
\phantomsection\label{concept:mqa}
\phantomsection\label{concept:gqa}
\phantomsection\label{partq:architecture}

标准 \conceptlink{multihead}{多头注意力}（MHA）可以为每个 Query 头保存自己的 Key/Value 头。
在自回归生成中，过去 Token 的 K/V 需要反复使用，于是产生 KV Cache。
\textbf{MQA} 让多个 Query 头共享一组 K/V；
\textbf{GQA} 让每组 Query 头共享一组 K/V，
是在“完全独立”和“全部共享”之间的一种折中。
\par\medskip
\tagstructbegin{tag=Figure,alttext={四个蓝色 Query 头与 Key Value 头的连接：MHA 一对一共四组 KV；GQA 每两个 Query 共享一组共两组 KV；MQA 四个 Query 全部共享一组 KV。灰色连线指示共享，不表示相同的查询内容。}}
\begin{center}
\begin{tikzpicture}[>=Latex,
  queryhead/.style={draw=blue,fill=blueback,rounded corners=2pt,minimum width=9mm,minimum height=7mm,font=\scriptsize},
  kvhead/.style={draw=teal,fill=greenback,rounded corners=2pt,minimum width=12mm,minimum height=8mm,font=\scriptsize},
  share/.style={->,draw=muted,line width=.55pt}]
\foreach \offset/\name/\label in {0/MHA/4组 KV,5.1/GQA/2组 KV,10.2/MQA/1组 KV}{
  \begin{scope}[xshift=\offset cm]
  \node[font=\sffamily\small\bfseries,text=ink] at (1.65,.72) {\name};
  \node[font=\scriptsize,text=muted] at (1.65,-2.07) {\label};
  \foreach \i in {0,1,2,3}{
    \node[queryhead] (q\name\i) at (.3+0.9*\i,0) {Q\the\numexpr\i+1\relax};
  }
  \end{scope}
}
\foreach \i in {0,1,2,3}{
 \node[kvhead] (kM\i) at (.3+0.9*\i,-1.23) {\textcolor{teal}{K}/\textcolor{orange}{V}};
 \draw[share] (qMHA\i.south)--(kM\i.north);
}
\foreach \i/\x in {0/.75,1/2.55}{
 \node[kvhead] (kG\i) at (5.1+\x,-1.23) {\textcolor{teal}{K}/\textcolor{orange}{V}};
}
\foreach \i/\j in {0/0,1/0,2/1,3/1}{
 \draw[share] (qGQA\i.south)--(kG\j.north);
}
\node[kvhead] (kQ) at (11.85,-1.23) {\textcolor{teal}{K}/\textcolor{orange}{V}};
\foreach \i in {0,1,2,3}{
 \draw[share] (qMQA\i.south)--(kQ.north);
}
\end{tikzpicture}
\captionof{figure}{四个 Query 头的三种 K/V 共享方式：MHA 4 组、GQA 2 组、MQA 1 组；多头潜在注意力（MLA）是另一类潜在压缩思路}
\label{fig:mha-gqa-mqa}
\end{center}
\tagstructend


为了比较缓存的\textbf{理论元素个数}，假设单层：
上下文长度 $T=100$，每头维度 $d_k=8$，
有 8 个 Query 头。如果 MHA 有 8 个 K/V 头：
\[
\underbrace{2}_{K,V}\times
\underbrace{100}_{T}\times
\underbrace{8}_{H_{KV}}\times
\underbrace{8}_{d_k}
=\answerresult{12{,}800\ \text{个数}}.
\]
若采用 2 个 K/V 头的 GQA：
\[
2\times100\times2\times8=\answerresult{3{,}200\ \text{个数}}.
\]
在这个简化例子里，缓存元素数缩小为原来的四分之一。
要换算成实际字节，还必须乘以每个元素占用的字节数、
模型层数和批量大小，并考虑缓存布局、分配及实现开销。

\par\Needspace{75mm}
\tagstructbegin{tag=Figure,alttext={四个 Query 头读取两组 KV 缓存。Q1和Q2分别算自己的注意力分数，但都读取第一组历史K和V；Q3和Q4读取第二组。每组缓存含春天和花两个历史位置。四个Q头保留四路独立输出，只存两组历史KV，而不是复制四组。}}
\begin{center}
\begin{tikzpicture}[>=Latex,x=1mm,y=1mm,
 q/.style={draw=blue,fill=blueback,rounded corners=2pt,minimum width=23mm,minimum height=8mm,font=\small},
 cache/.style={draw=flowcolor,fill=paper,rounded corners=2pt,text width=53mm,minimum height=23mm,align=center,font=\small}]
\foreach \n/\x in {1/0,2/32,3/77,4/109}{
 \node[q] (q\n) at (\x,0) {Q\n};
}
\node[cache] (a) at (16,-25) {缓存组 A：$K_A,V_A$\\春天：$k_{A,1},v_{A,1}$\\花：$k_{A,2},v_{A,2}$};
\node[cache] (b) at (93,-25) {缓存组 B：$K_B,V_B$\\春天：$k_{B,1},v_{B,1}$\\花：$k_{B,2},v_{B,2}$};
\foreach \n/\group in {1/a,2/a,3/b,4/b}{\draw[->,thick,muted] (q\n.south)--(\group.north);}
\node[font=\small,align=center] at (16,-48) {$o_1=\softmax(q_1K_A^\top)V_A$\\$o_2=\softmax(q_2K_A^\top)V_A$};
\node[font=\small,align=center] at (93,-48) {$o_3=\softmax(q_3K_B^\top)V_B$\\$o_4=\softmax(q_4K_B^\top)V_B$};
\node[font=\footnotesize,text=muted] at (54.5,-64) {为突出共享关系，图中省略缩放与掩码；各 Q 头分别计算权重和输出};
\end{tikzpicture}
\captionof{figure}{GQA 共享的是历史 K/V：四个查询头仍保留四路输出，历史缓存只存两组}
\label{fig:gqa-cache-read}
\end{center}
\tagstructend


\textbf{多头潜在注意力}（Multi-head Latent Attention，MLA）
由 \href{https://arxiv.org/abs/2405.04434}{DeepSeek-V2（arXiv:\allowbreak 2405.04434）}提出，
把每个 Token 的 K/V 压缩成一条低秩潜在向量来缓存，走另一条路线；
GQA 则直接减少独立 K/V 头数。选型时要比较质量和实际缓存字节数。

\begin{example}[breakable]
\noindent\emph{真实配置：Qwen2.5-0.5B 的 32K 缓存是多少？}\par
它的公开配置为 $24$ 层、$14$ 个 Query 头、$2$ 个 KV 头，
隐藏宽度 $896$，所以每头宽度 $896/14=64$。
对单请求、每个 K/V 元素 $2$ 字节、完整保存所有位置的缓存，
\[
B_{\rm token}=2\text{(K/V)}\times24\times2\times64\times2
=12{,}288\text{ 字节}=12\text{ KiB}.
\]
把 32K 明确定义成 $32{,}768$ 个位置，缓存本体为
\[
B_{32K}=12{,}288\times32{,}768=402{,}653{,}184\text{ 字节},
\]
即 $384$ MiB，约 $0.403$ 十进制 GB。
保持其他条件不变，改成 $14$ 个 KV 头的 MHA，缓存为原来 $7$ 倍：
$2{,}818{,}572{,}288$ 字节，约 $2.819$ GB。
这是满长度 K/V 数据的估算，未计权重、分页碎片和运行时缓冲；
批量中的多条请求分别计入其有效缓存位置。
配置来源见附录 F；\texttt{scripts/cache\_budget\_demo.py} 从锁定配置复算。
\end{example}

\subsection{只替换一次 Norm：落回前面的三维数字}

三维贯通模型在 Attention 前的第一行是
$\vect x^\top=(1.30,0.05,-0.30)$。
此前使用固定 $\gamma=1,\beta=0$ 的 LayerNorm：
\[
\operatorname{LN}(\vect x^\top)\approx(1.3833,-0.4368,-0.9464).
\]
如果\textbf{仅为了比较数值}，在同一输入上改算
\[
\operatorname{RMSNorm}(\vect x^\top)
=\frac{\vect x^\top}{
\sqrt{(1.30^2+0.05^2+(-0.30)^2)/3+10^{-5}}}
\approx\answerresult{(1.6865,0.0649,-0.3892)}.
\]
两种归一化\textbf{同样输出 $1\times3$}，
但 LayerNorm 先减均值、RMSNorm 不减均值，
因此下一步 $Q=XW_Q$ 的输入数值会改变。
上面只是\textbf{单组件数值对比}；
真正替换架构后，要在相同训练设置和任务上检验收敛速度与输出质量。

\subsection{读完三项必读，怎样比较架构？}

\begin{center}
\small
\renewcommand{\arraystretch}{1.3}
\tagpdfsetup{table/header-rows={1}}
\begin{tabularx}{\linewidth}{@{}>{\raggedright\arraybackslash}p{30mm}>{\raggedright\arraybackslash}p{45mm}>{\raggedright\arraybackslash}X@{}}
\toprule
设计部件 & 你应追问的输入/输出 & 首先检查的权衡\\
\midrule
位置表示 & Q/K 如何感知相对位置？ & 长度外推、位置泛化\\
归一化/残差 & 特征尺度和梯度怎样流动？ & 稳定性、训练成本\\
逐位置 FFN & 中间层扩到多宽？如何激活？ & 参数量、计算量、表示能力\\
MHA/MQA/GQA & K/V 头有多少？ & 缓存、质量、吞吐\\
注意力模式 & 哪些 Token 彼此可见？ & 长程依赖、复杂度\\
\bottomrule
\end{tabularx}
\end{center}

\subsection{［复算］从公开配置文件对照两种模型}

只凭宣传名称判断“用了哪些技术”容易出错。
\texttt{scripts/config\_table.py} 读取带来源和固定修订号的
\texttt{config.json} 字段快照，生成以下最小对照：

\begin{center}\small
\tagpdfsetup{table/header-rows={1}}
\begin{tabularx}{\linewidth}{@{}p{29mm}p{16mm}p{15mm}p{18mm}X@{}}
\toprule
参考模型 & 隐藏宽 & 层数 & Q/KV 头 & 位置／归一化\\\midrule
GPT-2 small & 768 & 12 & 12/12 & 学习式位置，LayerNorm\\
Qwen2.5-0.5B & 896 & 24 & 14/2 & RoPE，RMSNorm\\
\bottomrule
\end{tabularx}
\end{center}
后一配置写明 SiLU 与 \texttt{use\_sliding\_window=false}：此配置没有启用滑动窗口。

配置字段说明结构元数据；门控 FFN 的完整形式还要核对模型实现。


\textbf{复算入口：}
\texttt{python3 scripts/config\_table.py}
（只读仓库中固定版本的公开配置摘要，不下载权重）；
加入 \texttt{--verify-online} 才会联网核对原始 \texttt{config.json}。
来源：
\href{https://huggingface.co/openai-community/gpt2/blob/607a30d783dfa663caf39e06633721c8d4cfcd7e/config.json}{GPT-2 config}、
\href{https://huggingface.co/Qwen/Qwen2.5-0.5B/blob/060db6499f32faf8b98477b0a26969ef7d8b9987/config.json}{Qwen2.5 config}。
这张表只比较\textbf{结构元数据}；
模型质量与推理速度需由独立基准测量。

\subsection*{\texorpdfstring{{\symboliconfont ★}}{★} 选读：FFN、稀疏与线性状态的变体}
\noindent\textcolor{muted}{以下按需选读；首次阅读可直接跳到章末自测。}\par
\subsection{\texorpdfstring{{\symboliconfont ★}}{★} 选读：SwiGLU 与 MoE 怎样改变 FFN？}
\label{elective:architecture-extras}
\phantomsection\label{concept:router}
\phantomsection\label{concept:swiglu}
\phantomsection\label{concept:moe}

常规 \conceptlink{ffn}{FFN} 使用两次矩阵投影夹一个激活函数。
除了对特征做非线性变换，能否还根据当前输入，分别调节各项特征的传递强度？
例如一条分支给出内容 $(2,3)$，另一条给出调节系数 $(0.5,2)$，
对应位置相乘便得到 $(1,6)$：第一项减弱，第二项增强。
这种由另一条分支调节内容的做法叫作\textbf{门控}，两个分支都由输入和可训练权重产生。

SwiGLU 是一种门控 FFN：门分支使用 SiLU 激活，内容分支使用线性投影，相乘后再降回输出宽度：
\begin{equation}
\operatorname{FFN}(\vect x)
=\left[
\operatorname{SiLU}(\vect x^\top W_g)
\odot(\vect x^\top W_u)\right]W_{\rm down},
\label{eq:swiglu-illustration}
\end{equation}
其中 $\odot$ 是\textbf{对应位置逐元素相乘}，
$\operatorname{SiLU}(z)=z/(1+e^{-z})$。
其中 $W_g,W_u$ 将特征扩到同一中间宽度，$W_{\rm down}$ 再缩回。
门值是调节系数，可以大于 $1$ 或为负，并非只能开或关的 $0/1$ 开关。

另一个问题是：希望增加 FFN 的总容量，却又希望控制每个 Token 的计算量，怎么办？
可以准备多套分别带参数的 FFN，每套叫一个\textbf{专家}；每次只选其中一部分处理当前 Token。
这种组织方式叫作\textbf{混合专家（MoE）}，负责根据输入打分、挑选专家的模块叫\textbf{路由器}。

一个最简单的路由示例：三个专家的得分是 $(2,1,0)$，
如果只选择最高的一个专家，本次仅计算专家 1，
其他专家并未被本次 Token 调用。
选择 Top-2 时计算专家 1、2，再对入选的两个分数做 Softmax 重归一化：
\[
w_1=\frac{e^2}{e^2+e^1}\approx\boxed{0.731},\qquad
w_2=\frac{e^1}{e^2+e^1}\approx\boxed{0.269},\qquad
\text{输出}=0.731\,f_1(\vect x)+0.269\,f_2(\vect x).
\]
先对三个分数做 Softmax 得 $(0.665,0.245,0.090)$，再把前两项除以它们的和 $0.910$，结果相同。
实际系统还需要处理负载均衡、通信与容量限制。

\par\Needspace{15\baselineskip}
\tagstructbegin{tag=Figure,alttext={MoE 的单 Token 路由流程。蓝色 Token 进入路由器，三个专家的分数分别为二、一、零；Top-1 只将 Token 送给专家一，其他两个专家此次不计算，专家一的输出送往汇合位置。}}
\begin{center}
\begin{tikzpicture}[>=Latex,
 block/.style={draw=linegray,fill=paper,rounded corners=2pt,
   minimum width=20mm,minimum height=9mm,align=center,font=\sffamily\small},
 expert/.style={draw=flowcolor,fill=linegray!15,rounded corners=2pt,
   minimum width=24mm,minimum height=9mm,align=center,font=\sffamily\small}]
\node[block] (token) at (0,0) {Token};
\node[block] (router) at (3.0,0) {Router\\$(2,1,0)$};
\node[expert] (one) at (6.8,1.15) {专家 1\\Top-1 选中};
\node[block,fill=paper] (two) at (6.8,0) {专家 2\\未选中};
\node[block,fill=paper] (three) at (6.8,-1.15) {专家 3\\未选中};
\node[block] (output) at (10.3,1.15) {输出};
\draw[->,thick,flowcolor] (token)--(router);
\draw[->,thick,flowcolor] (router.east)--(one.west);
\draw[->,thick,flowcolor] (one)--(output);
\draw[densely dashed,muted] (router.east)--(two.west);
\draw[densely dashed,muted] (router.east)--(three.west);
\end{tikzpicture}
\captionof{figure}{单 Token 的 MoE Top-1 路由：实线是本次实际计算路径，虚线代表未选中的专家；换成 Top-2 会激活两位专家}
\label{fig:moe-token-routing}
\end{center}
\tagstructend

\begin{warningbox}
MoE 的\emph{总参数量}与每 Token 激活量分开计数。
每个 Token 只激活部分专家，设备仍需存储或调度专家权重，并承担路由与通信成本。
\end{warningbox}

\subsection{\texorpdfstring{{\symboliconfont ★}}{★} 选读：长上下文与稀疏访问}
\phantomsection\label{concept:sliding-window-attention}
\phantomsection\label{concept:linear-attention}
\phantomsection\label{concept:sparse-attention}

对长度为 $T$ 的普通全量 Attention，分数矩阵有 $T^2$ 个位置；
如果 $T$ 从 $100$ 增到 $200$，潜在分数位置从 $10{,}000$
变成 $40{,}000$。这解释了长序列计算面临的压力。

\begin{center}
\small
\renewcommand{\arraystretch}{1.3}
\tagpdfsetup{table/header-rows={1}}
\begin{tabularx}{\linewidth}{@{}>{\raggedright\arraybackslash}p{33mm}>{\raggedright\arraybackslash}X@{}}
\toprule
设计方向 & 为何这样改，以及保留什么代价\\
\midrule
局部/窗口 Attention & 仅计算附近的位置；远距离信息依靠跨层或额外全局通路传递\\
稀疏/选择性 Attention & 挑部分位置计算，减少无用比较；需设计选择策略\\
线性 Attention & 改变权重计算或状态更新形式，在特定条件下改善随长度的扩展性；按具体的核或状态更新公式定义其读取结果\\
混合结构 & 结合不同注意力或状态更新模块，按需求权衡质量、延迟和内存\\
\bottomrule
\end{tabularx}
\end{center}

\par\Needspace{16\baselineskip}
\tagstructbegin{tag=Figure,alttext={同样五个 Token 的两种因果注意力可见性。左侧全量因果注意力五行依次允许一、二、三、四、五个 Key，总计十五个位置对；右侧窗口大小二依次允许一、二、二、二、二个 Key，共九个位置对。蓝色方格表示允许访问，浅色方格表示不可见。}}
\begin{center}
\begin{tikzpicture}[x=.68cm,y=.68cm]
\foreach \shift/\kind/\label in {0/full/全量因果：15,6/local/窗口 2：9}{
  \begin{scope}[xshift=\shift cm]
  \node[font=\sffamily\small\bfseries,text=ink] at (2.5,.85) {\label};
  \foreach \r in {0,...,4}{
   \foreach \c in {0,...,4}{
    \ifnum\shift=0
      \pgfmathtruncatemacro{\active}{ifthenelse(\c<=\r,1,0)}
    \else
      \pgfmathtruncatemacro{\active}{ifthenelse((\c<=\r)&&(\c>=\r-1),1,0)}
    \fi
    \ifnum\active=1
      \fill[blue!65] (\c,-\r) rectangle ++(.84,-.84);
    \else
      \fill[paper] (\c,-\r) rectangle ++(.84,-.84);
    \fi
    \draw[linegray] (\c,-\r) rectangle ++(.84,-.84);
   }
  }
  \node[font=\scriptsize,text=muted] at (2.5,-5.25) {横轴 Key 位置，纵轴 Query 位置};
  \end{scope}
}
\end{tikzpicture}
\captionof{figure}{长度 $T=5$ 时，全量因果 Attention 与宽度 $w=2$ 的滑动窗口；深蓝格表示可访问，浅色格表示不可见，允许位置对由 15 减为 9}
\label{fig:causal-versus-window}
\end{center}
\tagstructend

判断扩展性时，先标明节约的是训练分数矩阵、解码 KV Cache，还是单 Token 计算。
这三个对象分别计量，后面的例子据此对照。

\subsection{\texorpdfstring{{\symboliconfont ★}}{★} 选读：线性 Attention 与 Delta 规则}

\textbf{将历史压入固定形状的状态，就能逐步写入与读取。}
沿用行向量约定，每来一对 Key--Value，就把外积 $k_i^\top v_i$（$d_k\times d_v$）加进状态：
\[
\underbrace{\mathcal{M}_t}_{d_k\times d_v}
=\sum_{i\le t}k_i^\top v_i,\qquad
\underbrace{y_t}_{1\times d_v}=q_t\mathcal{M}_t.
\]
查询时只读这个固定大小的状态，长度增加也不必保存每个位置的 K/V。

\begin{example}
\noindent\emph{两条关联如何汇成一个状态}\par
\localfixture{二维行向量状态，独立于三维 Transformer 贯通参数}
若 $k_1=(1,0)$、$v_1=(2,0)$，
$k_2=(0,1)$、$v_2=(0,3)$，
则外积相加得到 $\mathcal{M}_2=\begin{bmatrix}2&0\\0&3\end{bmatrix}$。
用 $q_2=(1,1)$ 查询，$q_2\mathcal{M}_2=(2,3)$；
两条 Value 的贡献直接相加，保留了累计幅度。
\end{example}

\phantomsection\label{sec:modern-delta-rule}
\textbf{Key 不正交时，直接相加会串扰。}
先写入 $k_A=(1,0)$、$v_A=(1,0)$，得 $\mathcal{M}_1=k_A^\top v_A=\begin{bmatrix}1&0\\0&0\end{bmatrix}$。
再来 $k_B=(0.6,0.8)$、$v_B=(0,1)$：若直接加上 $k_B^\top v_B$，
\[
\mathcal{M}_2^{\rm add}=\begin{bmatrix}1&0.6\\0&0.8\end{bmatrix},\qquad
k_B\mathcal{M}_2^{\rm add}=(0.6,\,1),
\]
多出的 $0.6$ 来自旧记忆 A，因为 $k_Bk_A^\top=0.6\ne0$。
\textbf{Delta 规则}先读旧状态，只写入还没记住的残差：
\[
r_t=v_t-k_t\mathcal{M}_{t-1},\qquad
\boxed{\mathcal{M}_t=\mathcal{M}_{t-1}+\eta\, k_t^\top r_t}.
\]
取 $\eta=1$：$r_B=(0,1)-(0.6,0)=(-0.6,1)$，
\[
\mathcal{M}_2=\mathcal{M}_1+\begin{bmatrix}0.6\\0.8\end{bmatrix}(-0.6,1)
=\begin{bmatrix}0.64&0.6\\-0.48&0.8\end{bmatrix},\qquad
k_B\mathcal{M}_2=\boxed{(0,1)}.
\]
B 的读出被精确修正。Gated DeltaNet 把这种残差写入与遗忘门控结合，用于高效序列建模。
写入 B 对旧 Key A 有什么影响、反复写入会收敛到哪里，
\hyperref[part:memory]{联想记忆部分}用外积视角再算一遍（可选延伸）。


\begin{boundarybox}
这里展示线性状态的最小算例。标准 Softmax Attention 还对每行权重归一化；
线性变体的特征映射、归一化与遗忘规则会改变输出。
Gated DeltaNet 的完整算法还要指定门控与参数化，本文只展示残差写入的动机。
\end{boundarybox}

\subsection{\texorpdfstring{{\symboliconfont ★}}{★} 选读索引：其他现代架构术语}

以 GPT-2 类 Decoder-only Transformer 为坐标系，按被替换的部件认识新设计。
Dense 使用相关稠密计算；稀疏 Attention 选择位置，稀疏 MoE 选择专家。

\begin{center}
\small
\renewcommand{\arraystretch}{1.25}
\tagpdfsetup{table/header-rows={1}}
\begin{tabularx}{\linewidth}{@{}>{\raggedright\arraybackslash}p{29mm}>{\raggedright\arraybackslash}p{41mm}>{\raggedright\arraybackslash}X@{}}
\toprule
经典部件（示例） & 可替换的方向 & 首先检查什么，保留什么代价\\\midrule
学习式绝对位置 & RoPE、YaRN、NoPE & 位置编码怎样进入 Q/K；扩展长度需长程实测\\
LayerNorm / Pre-Norm & RMSNorm 等 & 是否中心化、参数量、数值稳定性\\
缩放 $QK^\top$ & QK-Norm 等 & Q/K 是否归一化；分数尺度和 Softmax 饱和\\
Dense FFN（GELU） & SwiGLU / 稀疏 MoE & 门控、活跃专家、通信与总参数量\\
MHA 的 K/V 头 & MQA、GQA、MLA & 每层 KV Cache 布局、元素数与质量\\
简单残差旁路 & 超连接（HC）/ mHC & 多路径信息流与训练稳定性\\
Dense 全量 Attention & 局部/稀疏 Attention、Gated DeltaNet、Mamba 或混合层 & 长程访问、缓存、计算与训练/解码差别\\
\bottomrule
\end{tabularx}
\end{center}

这里 \textbf{YaRN} 是基于 RoPE 的长上下文扩展方向，
\textbf{NoPE} 表示某些设计不在该层显式加入位置编码，
顺序还可通过因果可见性等结构进入计算。
\textbf{QK-Norm} 尝试控制 Q/K 的尺度与分数饱和；
\textbf{超连接（Hyper-Connections, HC）}用多条可学习的残差通路替代单一残差旁路，
\textbf{mHC（Manifold-Constrained Hyper-Connections）}再约束这些通路的混合矩阵以恢复恒等映射性质；
这些设计要连同连接路径与约束一起理解。
可参考
\href{https://arxiv.org/abs/2309.00071}{YaRN}、
\href{https://aclanthology.org/2020.findings-emnlp.379/}{QK-Norm}、
\href{https://arxiv.org/abs/2409.19606}{HC}
与 \href{https://arxiv.org/abs/2512.24880}{mHC} 的原始论文。
\textbf{Mamba} 是选择性状态空间模型的一类路线，
\textbf{Gated DeltaNet} 采用门控 Delta 状态更新的高效序列建模路线；
两者按各自的状态更新公式建立比较。
参见 \href{https://arxiv.org/abs/2312.00752}{Mamba} 和
\href{https://arxiv.org/abs/2412.06464}{Gated DeltaNet}。

\subsection{\texorpdfstring{{\symboliconfont ★}}{★} 选读案例：因果编码器—解码器重新分工}
\phantomsection\label{elective:causal-encoder-decoder}

第~\ref*{intro:22-transformer-architecture}~章区分了原始 Encoder--Decoder 与 Decoder-only。
新架构也会重新组合编码与解码：截至 2026-10-25（UTC），
\href{https://huggingface.co/deepseek-ai/DeepSeek-V4.1-Flash}{DeepSeek-V4.1-Flash 官方模型卡}
介绍的\textbf{因果编码器—解码器（Causal Encoder--Decoder, CED）}
把 $40$ 层组织为 $20$ 层因果编码器加 $20$ 层解码器，整个模型仍按自回归方式生成文本；
解码器的全局 KV Cache 由编码器最后一层的隐藏状态投影得到。
\href{https://www.deepseek.com/en/news/deepseek-v4-1-flash/}{官方发布公告}给出的激活参数是：
处理输入时每 Token $8$B，生成输出时每 Token $16$B（MoE 骨干共 $552$B 参数）。
原始翻译模型的编码器双向读取整句源文；这里的编码器冠以“因果”，读结构图时先核对它的可见范围。
按本章的问法，它替换的是“每层同时负责读入与生成”的统一堆叠；
首先检查 KV Cache 来自哪一层、预填充与解码各激活多少参数。
这些是厂商公布的结构描述，质量与速度需按本章末的条件另行对照。

\subsection{\texorpdfstring{{\symboliconfont ★}}{★} 选读复算：MLA 缓存与窗口}

\localfixture{假设参数下的单层缓存与五位置窗口；不对应特定发布模型}
取 $T=8$、4 个 Query 头、每头 Key/Value 宽度 $2$。
普通 MHA 的 Key/Value 缓存元素数为
$2\times8\times4\times2=\answerresult{128}$。
若 GQA 只保留 2 个 K/V 头：
$2\times8\times2\times2=\answerresult{64}$。

再假设一种 MLA 式压缩设计，每个 Token 缓存
宽度为 $3$ 的共享潜在向量和宽度为 $2$ 的位置相关部分，
则\textbf{这份假设}的缓存数为
$8\times(3+2)=\answerresult{40}$。


对五个 Token，用窗口 $w=3$ 的\textbf{包含当前位置}的因果掩码，
行是 Query 位置，列是 Key 位置；1 表示允许读取：
\[
\boxed{\begin{bmatrix}
1&0&0&0&0\\
1&1&0&0&0\\
1&1&1&0&0\\
0&1&1&1&0\\
0&0&1&1&1
\end{bmatrix}}.
\]
最后一个 Token 只能读取位置 3、4、5（这里位置从 1 编号）。
窗口掩码有 $1+2+3+3+3=\answerresult{12}$ 个可见位置对，
全量因果 Mask 则有 $15$ 个。
\begin{boundarybox}
\textbf{计算口径：}以上 $128/64/40$ 是\textbf{示例结构的缓存元素数}；
实际 MLA 还涉及位置分支、张量布局、投影与解码重建。
掩码的 $12/15$ 是\textbf{可见位置对数}；
实际显存和运行时间还取决于实现、数据类型、Kernel 与设备。
\end{boundarybox}

\noindent{\sffamily\footnotesize\color{blue}\textbf{［复算］}}
运行 \texttt{python3 scripts/architecture\_demo.py}
可以检查 RoPE 平移例、MoE Top-2 权重、缓存元素数与 $5\times5$ 掩码的 12 个可见位置；
\texttt{python3 scripts/memory\_demo.py} 复算上文的串扰与 Delta 写入。


\begin{boundarybox}[title={适用条件：结构计数与实际性能}]
本章的缓存字节数、激活专家数与位置对数按给定配置计算。
性能对照统一输入长度、批量、精度、硬件和目标质量，另测长程检索、吞吐与延迟；
计数节省先解释瓶颈，实际收益由实现与负载决定。
\end{boundarybox}

\chapterchecktasks{22d-modern-transformers}{理解与复算}{
\item \textbf{比较：}RoPE 与 RMSNorm 分别主要改动位置表示还是归一化？RMSNorm 是否和 LayerNorm 一样先减去均值？\chapterinlinehint{分别对照位置旋转与归一化公式。}
\item \textbf{手算：}保持层数与头维不变，MHA 有 $12$ 个 K/V 头，GQA 只有 $3$ 个，单 Token 的 K/V 缓存元素比是多少？如果长度翻倍，相对原 MHA、原长度的缓存规模是多少？\chapterinlinehint{只保留缓存公式中的头数与长度变化项。}
\item \textbf{找错：}团队说“KV 元素数减少四分之三，所以端到端生成必定快四倍”。至少列出两个尚未控制的影响因素。\chapterinlinehint{列出缓存之外的运算、带宽与调度成本。}
}


\par\smallskip
现代 Transformer 在位置表示、数值尺度、参数使用方式和上下文访问方式上提供了不同选择。
判断一项技术，先问它解决了什么已知问题，在哪种计算或数据条件下有收益，付出了什么代价。
下一章从这些具体选择中抽离出来，讨论怎样用实验、来源与适用条件评价任何新的模型架构或 Agent 技术。
